\chapter*{Streszczenie}

%Krótki wstęp.

%Część główna, która powinna być nieco dłuższa. Całe streszcznie powinno
%zająć około pół strony.

%Krótkie podsumowanie wniosków, wyników i ewentualnych proponowanych
%kolejnych kroków.

Silniki cieplne to maszyny konwertujące dostarczane do nich ciepło na pracę. W silnikach klasycznych czynnikiem roboczym najczęściej jest gaz, który czasem może zmienić stan na ciekły. W silnikach kwantowych natowmiast czynnikiem roboczym są układu kwantowe. Podczas pracy mogą być one w kontakcie ze zbiornikami ciepła o różnych temperaturach oraz w kontrolowanych warunkach zewnętrznych. Zmiana energii związana z samorzutną zmianą obsadzenia stanów kwantowych interpretowana jest jako ciepło, a jej zmiana poprzez kontrolowany parametr zewnętrzny jako praca. Silniki kwantowe dla pewnych cykli i czynników roboczych wykazują wyższe sprawności niż ich klasyczne odpowiedniki.

Kwantowe silniki cieplne są rezultatem połączenia analizy klasycznych silników cieplnych z układami kwantowymi jako czynnikiem roboczym. Pierwsze prace nad nimi rozpoczęły się od masera. Jest to układ kwantowy o przynajmniej trzech poziomach energetycznych. Jest on wzbudzany do stanu najwyższego, po czym następują dwa przejścia na kolejne stany niższe. Jedno z nich jest emisją wymuszoną, która jest analogiem wykonanej pracy.

W pracy przedstawiłem wyprowadzenie równań opisujących wybrane cykle termodynamiczne pracujące z wykorzystaniem dwóch układów kwantowych w roli czynnika roboczego: oscylatora harmonicznego oraz układu o spinie \(\inv{2}\). Wykonałem obliczenia numeryczne w celu znalezienia sprawności w punkcie maksymalnej mocy.

% Define 5-10 keywords which can describe what you're writing about
\bigskip
\noindent
\textbf{Słowa kluczowe:} silniki cieplne, termodynamika kwantowa

% Choose the field from the OECD list:
% https://en.wikipedia.org/wiki/Fields_of_Science_and_Technology
% With some research, you should find the whole list officially
% traslated into Polish.
% The below should be fine for PG ETI (informatyka)
\bigskip
\noindent
\textbf{Dziedzina nauki i techniki, zgodnie z wymogami OECD:}
nauki fizyczne
